%%%%%%%%%%%%
%
% $Autor: C. Crety, J. Perewersenko, H. Mey $
% $Datum: 2025-07-10  $
% $Pfad: PushButton $
%
%%%%%%%%%%%%

% Structure
\chapter{Taster}\index{Push Button! Built-in Push Button}

\section{Allgemeine Beschreibung des Tasters}

Ein Taster ist ein einfacher Schaltmechanismus, der zur Steuerung verschiedener Geräte und Prozesse verwendet wird. Er besteht typischerweise aus harten Materialien wie Kunststoff oder Metall.
Die Oberfläche eines Tasters ist so gestaltet, dass sie leicht mit dem Finger oder der Hand gedrückt werden kann. Beim Drücken des Tasters wird ein elektrischer Stromkreis entweder geschlossen oder geöffnet.
In industriellen und kommerziellen Anwendungen können Taster miteinander verbunden werden, sodass das Drücken eines Tasters einen anderen freigibt. Not-Aus-Taster, oft mit großen, pilzförmigen Köpfen, erhöhen die Sicherheit von Maschinen und Anlagen. Um Aufmerksamkeit zu erzeugen und Rückmeldung beim Betätigen zu geben, werden Taster manchmal mit Signalleuchten kombiniert. Eine Farbcodierung ist üblich, um Taster bestimmten Funktionen zuzuordnen (z. B. Rot für Stopp, Grün für Start). \cite{tietze2021elektronik}


\section{Spezifische Beschreibung des Tasters}

Im Schrittmotor-Demonstrator wird ein Taster als zentrale Bedieneinheit zur Steuerung der Spindelrotation eingesetzt. Dabei erfüllt er zwei wesentliche Funktionen: das Starten bzw. Fortsetzen der Spindelbewegung sowie das Stoppen der Rotation. Zur eindeutigen Unterscheidung der beiden Betriebszustände kommen zwei farblich codierte Taster zum Einsatz – Grün für „Start/Weiterlauf“ und Rot für „Stopp“.
Durch das Drücken des grünen Tasters wird ein entsprechendes Signal an den Mikrocontroller übermittelt, woraufhin der Motor aktiviert oder – nach einem vorherigen Stopp – fortgesetzt wird. Der rote Taster hingegen unterbricht den Antrieb sofort, indem er ein Stoppsignal an den Mikrocontroller sendet.


\subsection{Anschluss des Tasters mit dem Arduino Nano 33 BLE Sense}

Der Arduino Nano 33 BLE Sense verfügt über 18 digitale Eingänge, die zur Erfassung von Schaltzuständen verwendet werden können. Die Taster werden jeweils an einen digitalen Eingang angeschlossen. Beim Drücken eines Tasters wird der entsprechende Eingang auf LOW (GND) gezogen, während bei geöffnetem Zustand ein interner Pullup-Widerstand den Pegel auf HIGH hält. Diese Konfiguration verhindert ungewollte Schaltzustände durch Störungen


\section{Test}

\subsection{Code-Beispiel}


\begin{lstlisting}
	// Pin definitions
	const int startButtonPin = 2;  // Green button: Start
	const int stopButtonPin  = 3;  // Red button: Stop
	
	void setup() 
	{
		// Start serial communication
		Serial.begin(9600);
		
		// Set pins as inputs with internal pull-up resistors
		pinMode(startButtonPin, INPUT_PULLUP);
		pinMode(stopButtonPin, INPUT_PULLUP);
	}
	
	void loop() 
	{
		// Read button states (LOW means pressed due to pull-up)
		bool startPressed = digitalRead(startButtonPin) == LOW;
		bool stopPressed  = digitalRead(stopButtonPin) == LOW;
		
		// Print status when button is pressed
		if (startPressed) 
		{
			Serial.println("START button pressed");
			delay(300);  // Basic debounce delay
		}
		
		if (stopPressed) 
		{
			Serial.println("STOP button pressed");
			delay(300);  // Basic debounce delay
		}
	}
	
\end{lstlisting}



